% ============================================================
% 第5章 化学反应与电能（选必1第四章，课件未覆盖，补充）
% ============================================================
\chapter{化学反应与电能}

\begin{daoxue}
\textbf{知识主线：}自发的氧化还原反应（氧化、还原分区进行）$\to$ 原电池（化学能$\to$电能）$\to$ 化学电源；外加电源（强制发生氧化还原）$\to$ 电解池（电能$\to$化学能）$\to$ 氯碱、电镀、精炼、冶金；两者结合 $\to$ 金属的腐蚀与防护。核心思想只有一条：\textbf{负极（阳极）失电子发生氧化反应，正极（阴极）得电子发生还原反应；电子守恒}。

\textbf{高考考情：}广东卷每年选择题中有\textbf{两道}电化学题：2024年第5题（牺牲阳极）、第16题（离子交换膜电解与计算）；2025年第14题（水系可充电电池）、第16题（燃料电池催化机理）；2026年第6题（海上风电防腐）、第16题（电解还原 \ce{CO2} 与海水脱盐）。大题中也常出现电镀、电解回收（2025年第18题、2026年第17、18题）。

\textbf{学习目标：}①判断电极、电子与离子的流向；②规范书写各种介质中的电极反应式；③会分析离子交换膜的作用；④熟练进行电子守恒计算；⑤理解腐蚀原理与防护方法。
\end{daoxue}

\section{原电池的工作原理}

\subsection{从氧化还原到电流}
把锌片直接放入 \ce{CuSO4} 溶液，\ce{Zn + Cu^2+ = Zn^2+ + Cu}，电子直接转移，化学能变成热能。若让氧化反应和还原反应\textbf{在两个区域分别进行}，电子只能通过外电路导线从锌流向铜，就形成了电流——这就是原电池。

\begin{center}
\begin{tikzpicture}[scale=0.8,font=\small]
\draw[thick] (0,0) rectangle (2.6,2.2); \draw[thick] (4.4,0) rectangle (7,2.2);
\fill[mainblue!10] (0.02,0.02) rectangle (2.58,1.6); \fill[mainblue!18] (4.42,0.02) rectangle (6.98,1.6);
\draw[very thick,gray] (1.0,0.4)--(1.0,3.0); \node at (0.6,2.5) {Zn}; \draw[very thick,orange!70!black] (6.0,0.4)--(6.0,3.0); \node at (6.4,2.5) {Cu};
\draw[thick] (1.0,3.0)--(1.0,3.6)--(6.0,3.6)--(6.0,3.0);
\draw[-Stealth,warnred,thick] (2.6,3.75)--(4.4,3.75) node[midway,above]{\scriptsize 电子 $e^-$};
\draw[thick,deepgreen] (1.9,1.3)--(1.9,2.6)--(5.1,2.6)--(5.1,1.3); \draw[thick,deepgreen] (2.2,1.3)--(2.2,2.3)--(4.8,2.3)--(4.8,1.3);
\node[deepgreen] at (3.5,2.95) {\scriptsize 盐桥（KCl 琼胶）};
\node at (1.3,0.8) {\scriptsize \ce{ZnSO4}}; \node at (5.6,0.8) {\scriptsize \ce{CuSO4}};
\node[align=center] at (1.3,-0.6) {\scriptsize 负极：\ce{Zn - 2e- = Zn^2+}};
\node[align=center] at (5.7,-0.6) {\scriptsize 正极：\ce{Cu^2+ + 2e- = Cu}};
\node[align=left,anchor=west] at (7.4,1.6) {\scriptsize 盐桥中 \ce{K+} 移向正极区，\\[-1pt]\scriptsize \ce{Cl-} 移向负极区，\\[-1pt]\scriptsize 使两溶液保持电中性};
\end{tikzpicture}
\end{center}

\begin{gainian}[原电池]
将化学能转化为电能的装置。构成条件：①有自发进行的氧化还原反应；②两个电极；③电解质溶液（或熔融电解质）；④形成闭合回路。
\textbf{负极}：电子流出，发生\textbf{氧化反应}；\textbf{正极}：电子流入，发生\textbf{还原反应}。外电路电子由负极流向正极；内电路阳离子移向正极，阴离子移向负极。
\end{gainian}

\textbf{盐桥的作用：}①连通内电路；②平衡两池电荷（否则 \ce{ZnSO4} 溶液阳离子过多、\ce{CuSO4} 溶液阴离子过多，反应很快停止）；③使氧化剂与还原剂不直接接触，电流稳定、能量利用率高。用离子交换膜代替盐桥也是同样的道理。

\begin{yicuo}
负极不一定是“更活泼的金属”，要看谁与电解质溶液发生自发反应：Mg、Al 插入 NaOH 溶液，Al 与 NaOH 反应，\textbf{Al 作负极}（\ce{Al - 3e- + 4OH- = [Al(OH)4]-}）；Fe、Cu 插入浓硝酸，Fe 钝化，\textbf{Cu 作负极}。另外，电子只在导线中流动，\textbf{“电子不下水，离子不上岸”}。
\end{yicuo}

\subsection{电极反应式的书写}
\begin{fangfa}[三步写电极反应]
①\textbf{写总反应、找得失}：标出化合价变化，确定负极反应物、正极反应物及转移电子数。②\textbf{按介质补粒子}：酸性介质用 \ce{H+} 和 \ce{H2O} 配平，碱性介质用 \ce{OH-} 和 \ce{H2O}，熔融碳酸盐用 \ce{CO3^2-}，固体氧化物用 \ce{O^2-}；酸性介质中不能出现 \ce{OH-}，碱性介质中不能出现 \ce{H+}，也不能生成 \ce{CO2}（应生成 \ce{CO3^2-}）。③\textbf{查守恒}：原子、电荷、电子守恒；两极反应式相加（电子数相同）应得总反应。
\end{fangfa}

\begin{liti}[不同介质中燃料电池的电极反应]
写出下列燃料电池的电极反应式：(1) 甲醇—空气，质子交换膜（酸性）；(2) 肼（\ce{N2H4}）—空气，KOH 溶液；(3) CO—空气，熔融 \ce{K2CO3}（正极通入空气和 \ce{CO2}）；(4) \ce{CH4}—空气，固体氧化物电解质（能传导 \ce{O^2-}）。
\tcblower
\silu 先写正极（只由介质决定），再用“总反应$-$正极反应”得负极。\\
\jie (1) 正：\ce{O2 + 4e- + 4H+ = 2H2O}；负：\ce{CH3OH - 6e- + H2O = CO2 ^ + 6H+}。\\
(2) 正：\ce{O2 + 4e- + 2H2O = 4OH-}；负：\ce{N2H4 - 4e- + 4OH- = N2 ^ + 4H2O}（总反应 \ce{N2H4 + O2 = N2 + 2H2O}）。\\
(3) 正：\ce{O2 + 4e- + 2CO2 = 2CO3^2-}；负：\ce{CO - 2e- + CO3^2- = 2CO2}。负极生成的 \ce{CO2} 可循环到正极使用。\\
(4) 正：\ce{O2 + 4e- = 2O^2-}；负：\ce{CH4 - 8e- + 4O^2- = CO2 + 2H2O}。\\
\xianjing 碱性介质中 C 元素最终以 \ce{CO3^2-} 存在；熔融碳酸盐中不能出现 \ce{OH-}、\ce{H+}；固体氧化物中不能出现 \ce{H2O} 参与传导。
\end{liti}

\textbf{原电池原理的应用：}加快反应速率（锌与稀硫酸反应时滴入少量 \ce{CuSO4} 溶液，Zn 置换出 Cu 形成微电池）；比较金属活动性（作负极者一般更活泼）；设计原电池（选一个自发的氧化还原反应，还原剂作负极或负极反应物，氧化剂溶液作电解质）。

\section{化学电源}

\begin{center}
\small
\begin{tabularx}{\linewidth}{p{2.2cm}p{4.6cm}X}
\toprule
电池 & 总反应 & 电极反应 \\
\midrule
碱性锌锰电池（一次） & \ce{Zn + 2MnO2 + 2H2O = 2MnO(OH) + Zn(OH)2} & 负：\ce{Zn - 2e- + 2OH- = Zn(OH)2}；正：\ce{MnO2 + e- + H2O = MnO(OH) + OH-} \\
铅酸蓄电池（二次） & \ce{Pb + PbO2 + 2H2SO4 <=>[\text{放电}][\text{充电}] 2PbSO4 + 2H2O} & 放电负：\ce{Pb - 2e- + SO4^2- = PbSO4}；放电正：\ce{PbO2 + 2e- + 4H+ + SO4^2- = PbSO4 + 2H2O}；放电时两极质量都增加，\ce{H2SO4} 浓度减小 \\
锂离子电池（二次） & \ce{Li_xC6 + Li_{1-$x$}CoO2 <=>[\text{放电}][\text{充电}] 6C + LiCoO2} & 放电负：\ce{Li_xC6 - $x$e- = $x$Li+ + 6C}；放电正：\ce{Li_{1-$x$}CoO2 + $x$Li+ + $x$e- = LiCoO2}；\ce{Li+} 在两极间“摇摆” \\
氢氧燃料电池 & \ce{2H2 + O2 = 2H2O} & 酸性：负 \ce{H2 - 2e- = 2H+}，正 \ce{O2 + 4e- + 4H+ = 2H2O}；碱性：负 \ce{H2 - 2e- + 2OH- = 2H2O}，正 \ce{O2 + 4e- + 2H2O = 4OH-} \\
甲烷燃料电池（KOH） & \ce{CH4 + 2O2 + 2OH- = CO3^2- + 3H2O} & 负：\ce{CH4 - 8e- + 10OH- = CO3^2- + 7H2O}；正：\ce{2O2 + 8e- + 4H2O = 8OH-}；溶液碱性减弱 \\
\bottomrule
\end{tabularx}
\end{center}

\begin{benzhi}[燃料电池中燃料一定在负极，为什么？]
燃料（\ce{H2}、\ce{CH4}、\ce{CH3OH}、\ce{N2H4} 等）是还原剂，在负极失电子；\ce{O2} 是氧化剂，在正极得电子。正极反应只取决于介质：酸性 \ce{O2 + 4e- + 4H+ = 2H2O}；碱性或中性 \ce{O2 + 4e- + 2H2O = 4OH-}；熔融碳酸盐 \ce{O2 + 4e- + 2CO2 = 2CO3^2-}；固体氧化物 \ce{O2 + 4e- = 2O^2-}。写负极反应时，常用“总反应 $-$ 正极反应”得到，最不易出错。
\end{benzhi}

\textbf{二次电池的充电：}充电是放电的逆过程。\textbf{“负接负、正接正”}——电池负极接外电源负极（充电时作\textbf{阴极}），正极接外电源正极（作\textbf{阳极}）。充电时阴极反应是放电负极反应的逆反应，阳极反应是放电正极反应的逆反应。

\begin{fangfa}[新型电池题的解题步骤]
①\textbf{找变价}：由元素化合价变化（或“放电时某物质减少”）确定正、负极；②\textbf{定流向}：电子由负极经导线到正极，阳离子向正极迁移（充电时向阴极迁移）；③\textbf{写反应}：注意介质和膜允许通过的离子；④\textbf{算数量}：电子守恒，1 mol 电子对应电量 96500 C，即 26.8 A·h。
\end{fangfa}

\begin{liti}[全钒液流电池]
全钒液流电池的两极电解液分别含钒的不同价态离子，用质子交换膜隔开，总反应为
\[\ce{VO2+ + V^2+ + 2H+ <=>[\text{放电}][\text{充电}] VO^2+ + V^3+ + H2O}\]
（颜色：\ce{V^2+} 紫色，\ce{V^3+} 绿色，\ce{VO^2+} 蓝色，\ce{VO2+} 黄色）。(1) 写出放电时正、负极的电极反应式。(2) 放电时 \ce{H+} 通过质子交换膜向哪一极迁移？(3) 充电时阳极室溶液颜色如何变化？(4) 放电过程中每转移 1 mol 电子，正极室溶液中 $n(\ce{H+})$ 净变化多少？
\tcblower
\jie (1) V 元素：\ce{VO2+}（$+5$）$\to$\ce{VO^2+}（$+4$），被还原，在正极：\ce{VO2+ + 2H+ + e- = VO^2+ + H2O}；\ce{V^2+}（$+2$）$\to$\ce{V^3+}（$+3$），在负极：\ce{V^2+ - e- = V^3+}。(2) 移向正极。(3) 充电时阳极（放电正极）发生 \ce{VO^2+ - e- + H2O = VO2+ + 2H+}，溶液由蓝色变为黄色。(4) 每转移 1 mol 电子，正极反应消耗 2 mol \ce{H+}，同时有 1 mol \ce{H+} 经膜从负极室迁入，正极室 $n(\ce{H+})$ 净减少 1 mol。\\
\dianbo 分析离子交换膜问题时，“电极反应消耗或生成的离子”与“经膜迁移的离子”要分开算，再求净变化。
\end{liti}

\begin{zhenti}{2025年广东卷第14题\quad 来源：仓库所附真题文档}
一种高容量水系电池：电极Ⅰ为 S/MnS，电极Ⅱ为 \ce{MnO2}，两极室均为 \ce{MnSO4} 溶液，中间为离子交换膜。已知放电时电极Ⅱ上 \ce{MnO2} 减少；电极材料每转移 1 mol 电子，对应的理论容量为 26.8 A·h。下列说法错误的是
\xxb{充电时Ⅱ为阳极}{放电时Ⅱ极室中溶液的 pH 降低}{放电时负极反应为 \ce{MnS - 2e- = S + Mn^2+}}{充电时 16 g S 能提供的理论容量为 26.8 A·h}
\textbf{解析：}放电时 \ce{MnO2} 减少，被还原为 \ce{Mn^2+}，Ⅱ为正极：\ce{MnO2 + 2e- + 4H+ = Mn^2+ + 2H2O}，消耗 \ce{H+}，pH 升高，B 错误；充电时Ⅱ为阳极，A 对；Ⅰ为负极，C 对；充电时阴极 \ce{S + Mn^2+ + 2e- = MnS}，16 g S 为 0.5 mol，转移 1 mol 电子，D 对。答案：B。
\end{zhenti}

\section{电解池的工作原理}

\begin{gainian}[电解池]
借助电流引起氧化还原反应的装置（电能$\to$化学能）。与电源\textbf{负极}相连的是\textbf{阴极}，得电子发生还原反应；与电源\textbf{正极}相连的是\textbf{阳极}，失电子发生氧化反应。溶液中阳离子移向阴极，阴离子移向阳极。
\end{gainian}

\textbf{离子放电顺序（惰性电极，水溶液中）：}\\
阴极（阳离子得电子）：\ce{Ag+}$>$\allowbreak{}\ce{Fe^3+}$>$\allowbreak{}\ce{Cu^2+}$>$\allowbreak{}\ce{H+}（酸）$>$\allowbreak{}\ce{Pb^2+}$>$\allowbreak{}\ce{Sn^2+}$>$\allowbreak{}\ce{Fe^2+}$>$\allowbreak{}\ce{Zn^2+}$>$\allowbreak{}\ce{H+}（水）$>$\allowbreak{}\ce{Al^3+}$>$\allowbreak{}\ce{Mg^2+}$>$\allowbreak{}\ce{Na+}$>$\allowbreak{}\ce{Ca^2+}$>$\allowbreak{}\ce{K+}。\ce{Al^3+} 及其后的离子在水溶液中不放电，只能电解熔融物得到其单质。\\
阳极：\textbf{活性电极}（Pt、Au、石墨以外的金属）\textbf{优先自身失电子溶解}；惰性电极时阴离子放电：\ce{S^2-}$>$\allowbreak{}\ce{I-}$>$\allowbreak{}\ce{Br-}$>$\allowbreak{}\ce{Cl-}$>$\allowbreak{}\ce{OH-}$>$\allowbreak{}含氧酸根$>$\allowbreak{}\ce{F-}。

\begin{center}
\small
\begin{tabularx}{\linewidth}{p{2.2cm}p{2.6cm}Xp{2.6cm}}
\toprule
类型 & 实例 & 电极产物与溶液变化 & 复原方法 \\
\midrule
电解水型 & \ce{H2SO4}、NaOH、\ce{Na2SO4} & 阴极 \ce{H2}、阳极 \ce{O2}；溶质浓度增大 & 加水 \\
电解电解质型 & HCl、\ce{CuCl2} & 阴极 \ce{H2} 或 Cu，阳极 \ce{Cl2} & 通 HCl、加 \ce{CuCl2} \\
放 \ce{H2} 生碱型 & NaCl & 阴极 \ce{H2}，阳极 \ce{Cl2}；生成 NaOH & 通 HCl \\
放 \ce{O2} 生酸型 & \ce{CuSO4}、\ce{AgNO3} & 阴极 Cu（Ag），阳极 \ce{O2}；生成酸 & 加 CuO（\ce{Ag2O}） \\
\bottomrule
\end{tabularx}
\end{center}

\begin{yicuo}
\begin{enumerate}
\item 复原遵循“\textbf{出什么加什么}”：电解 \ce{CuSO4} 溶液析出 Cu 和 \ce{O2}，应加 CuO（或 \ce{CuCO3}），不能加 \ce{Cu(OH)2}（会多加水）——除非 \ce{Cu^2+} 放电完后又电解了水。
\item 水中 \ce{H+}、\ce{OH-} 放电时，电极反应可写作 \ce{2H2O + 2e- = H2 ^ + 2OH-}（阴极）、\ce{2H2O - 4e- = O2 ^ + 4H+}（阳极）；碱性溶液中阳极也可写 \ce{4OH- - 4e- = O2 ^ + 2H2O}，酸性溶液中阴极可写 \ce{2H+ + 2e- = H2 ^}，以溶液中大量存在的粒子为准。
\end{enumerate}
\end{yicuo}

\begin{liti}[混合溶液的分阶段电解]
用惰性电极电解 1 L 含 0.1 mol \ce{CuSO4} 和 0.1 mol NaCl 的混合溶液。当电路中通过 0.3 mol 电子时：(1) 分析电解的三个阶段；(2) 求阴、阳两极产生的气体的物质的量；(3) 电解后溶液呈酸性还是碱性？
\tcblower
\silu 按放电顺序“排队”：阴极 \ce{Cu^2+}$>$\ce{H+}；阳极 \ce{Cl-}$>$\ce{OH-}。分别算出每种离子放电所需电子数，再按电子数分段。\\
\jie (1) \ce{Cl-} 全部放电需 0.1 mol 电子，\ce{Cu^2+} 全部放电需 0.2 mol 电子。第一阶段（0～0.1 mol 电子）：阴极析 Cu，阳极放 \ce{Cl2}，相当于电解 \ce{CuCl2}；第二阶段（0.1～0.2 mol）：阴极析 Cu，阳极放 \ce{O2}，相当于电解 \ce{CuSO4}，生成 \ce{H2SO4}；第三阶段（0.2 mol 以后）：阴极放 \ce{H2}，阳极放 \ce{O2}，相当于电解水。\\
(2) 阳极：\ce{Cl2} 0.05 mol，\ce{O2} $\frac{0.3-0.1}{4}=0.05$ mol，共 0.1 mol；阴极：Cu 0.1 mol，\ce{H2} $\frac{0.3-0.2}{2}=0.05$ mol。\\
(3) 第二阶段生成 0.1 mol \ce{H+}（$2\ce{Cu^2+} + 2\ce{H2O} = 2\ce{Cu} + \ce{O2} + 4\ce{H+}$，0.1 mol 电子对应 0.1 mol \ce{H+}），第三阶段电解水不改变 \ce{H+} 的量，溶液呈酸性。
\end{liti}

\section{电解原理的应用}

\subsection{氯碱工业、电镀、精炼与电冶金}
\textbf{氯碱工业：}\ce{2NaCl + 2H2O ->[\text{电解}] 2NaOH + H2 ^ + Cl2 ^}。阳极（钛网）\ce{2Cl- - 2e- = Cl2 ^}；阴极（碳钢网）\ce{2H2O + 2e- = H2 ^ + 2OH-}。中间用\textbf{阳离子交换膜}，只允许 \ce{Na+} 由阳极室进入阴极室，既能得到较纯的 NaOH，又能防止 \ce{Cl2} 与 NaOH 反应、防止 \ce{H2} 与 \ce{Cl2} 混合爆炸。

\begin{center}
\small
\begin{tabularx}{\linewidth}{p{1.8cm}XXX}
\toprule
& 电镀（铁件镀铜） & 电解精炼铜 & 电冶金 \\
\midrule
阳极 & 镀层金属 Cu：\ce{Cu - 2e- = Cu^2+} & 粗铜（含 Zn、Fe、Ni、Ag、Au 等） & 惰性电极 \\
阴极 & 镀件（铁件） & 纯铜 & 析出金属 \\
电解质 & 含镀层金属离子的溶液（\ce{CuSO4}） & \ce{CuSO4} 溶液 & 熔融 NaCl、\ce{MgCl2}、\ce{Al2O3}（加冰晶石） \\
特点 & 电解质浓度基本不变 & Zn、Fe、Ni 先溶解，Ag、Au 形成阳极泥；$c(\ce{Cu^2+})$ 略减小 & \ce{2Al2O3 ->[\text{电解}][\text{冰晶石}] 4Al + 3O2 ^}（熔融）；不用 \ce{AlCl3}（共价化合物，熔融不导电） \\
\bottomrule
\end{tabularx}
\end{center}

\subsection{离子交换膜与多室电解}
\begin{fangfa}[判断离子交换膜种类的思路]
先根据电极反应确定各室“多出了什么电荷”（阳极室消耗阴离子或生成 \ce{H+}，阴极室消耗阳离子或生成 \ce{OH-}），再根据\textbf{产品要求}和\textbf{要阻止的副反应}判断哪种离子需要通过：阳离子交换膜只让阳离子通过，阴离子交换膜只让阴离子通过，质子交换膜只让 \ce{H+} 通过，双极膜在中间解离出 \ce{H+} 和 \ce{OH-} 分别移向阴、阳极。
\end{fangfa}

\begin{zhenti}{2024年广东卷第16题\quad 来源：仓库所附真题文档}
一种基于氯碱工艺的新型电解池可用于湿法冶铁：左侧电极室为 NaOH 溶液，电极上附着 \ce{Fe2O3}（电解后得到 Fe），中间为阳离子交换膜，右侧为饱和食盐水，电极上产生气体。下列说法\textbf{不正确}的是
\xxb{阳极反应：\ce{2Cl- - 2e- = Cl2 ^}}{阴极区溶液中 \ce{OH-} 浓度逐渐升高}{理论上每消耗 1 mol \ce{Fe2O3}，阳极室溶液减少 213 g}{理论上每消耗 1 mol \ce{Fe2O3}，阴极室物质最多增加 138 g}
\textbf{解析：}左侧为阴极：\ce{Fe2O3 + 6e- + 3H2O = 2Fe + 6OH-}，$c(\ce{OH-})$ 升高，B 对；右侧为阳极：\ce{2Cl- - 2e- = Cl2 ^}，A 对。每消耗 1 mol \ce{Fe2O3} 转移 6 mol 电子：阳极室放出 3 mol \ce{Cl2}（213 g），同时有 6 mol \ce{Na+}（138 g）迁移到阴极室，阳极室共减少 $213+138=351$ g，C 错；阴极室增加的是迁入的 \ce{Na+}，最多 138 g（\ce{Fe2O3} 转化为 Fe 和 \ce{OH-} 仍在该室），D 对。答案：C。
\end{zhenti}

\begin{liti}[双极膜电渗析]
双极膜由阳离子交换层和阴离子交换层复合而成，在直流电场作用下，膜间的水解离为 \ce{H+} 和 \ce{OH-}，分别向两侧迁移。用“阳极$|$双极膜$|$盐室$|$双极膜$|$阴极”结构处理 NaCl 溶液可同时得到 HCl 和 NaOH：盐室两侧分别为阴离子交换膜和阳离子交换膜。(1) 双极膜中解离出的 \ce{H+} 移向哪一极？(2) 盐室中 \ce{Na+}、\ce{Cl-} 分别向哪个方向迁移？在何处与 \ce{OH-}、\ce{H+} 结合？(3) 与直接电解 NaCl 溶液相比，这种方法的优点是什么？
\tcblower
\jie (1) \ce{H+} 是阳离子，移向阴极方向；\ce{OH-} 移向阳极方向。(2) \ce{Cl-} 透过阴离子交换膜向阳极方向迁移，与双极膜产生的 \ce{H+} 结合生成 HCl；\ce{Na+} 透过阳离子交换膜向阴极方向迁移，与双极膜产生的 \ce{OH-} 结合生成 NaOH。(3) 不在电极上产生 \ce{Cl2}、\ce{H2}（电极上仅发生少量电解水反应），能耗低，同时得到酸和碱，副产物少，更环保。
\end{liti}

\subsection{电化学计算}
核心是\textbf{电子守恒}：同一电路中各电极转移的电子数相等。常用关系：$4e^-\sim\ce{O2}\sim2\ce{H2}\sim2\ce{Cl2}\sim2\ce{Cu}\sim4\ce{Ag}\sim4\ce{H+}\sim4\ce{OH-}$。电量 $Q=It$，$n(e^-)=\dfrac{Q}{F}$（$F=96500$ C·mol$^{-1}$）。

\begin{liti}[电子守恒计算]
用惰性电极电解 500 mL \ce{CuSO4} 溶液，通电一段时间后阴极增重 0.64 g（无气体产生），设溶液体积不变。(1) 阳极产生气体的体积（标准状况）为多少？(2) 电解后溶液的 $c(\ce{H+})$ 为多少？(3) 若用 0.5 A 的恒定电流，需通电多少秒？
\tcblower
\jie (1) $n(\ce{Cu})=0.01$ mol，转移 0.02 mol 电子；阳极 \ce{2H2O - 4e- = O2 ^ + 4H+}，$n(\ce{O2})=0.005$ mol，$V=0.112$ L。(2) 生成 \ce{H+} 0.02 mol，$c(\ce{H+})=0.04$ mol·L$^{-1}$。(3) $t=\dfrac{nF}{I}=\dfrac{0.02\times96500}{0.5}=3860$ s。
\end{liti}

\section{金属的腐蚀与防护}

\textbf{化学腐蚀}：金属与接触的物质直接反应（如高温下铁与氯气）；\textbf{电化学腐蚀}：不纯的金属与电解质溶液接触形成原电池，较活泼的金属被氧化，更普遍、危害更大。

\begin{center}
\small
\begin{tabular}{lll}
\toprule
& 析氢腐蚀 & 吸氧腐蚀 \\
\midrule
条件 & 水膜酸性较强 & 水膜呈中性、弱酸性或碱性（最普遍） \\
负极 & \ce{Fe - 2e- = Fe^2+} & \ce{Fe - 2e- = Fe^2+} \\
正极 & \ce{2H+ + 2e- = H2 ^} & \ce{O2 + 4e- + 2H2O = 4OH-} \\
后续 & & \ce{Fe^2+ + 2OH- = Fe(OH)2}，\ce{4Fe(OH)2 + O2 + 2H2O = 4Fe(OH)3}，最终形成铁锈 \ce{Fe2O3*$x$H2O} \\
\bottomrule
\end{tabular}
\end{center}

\textbf{防护方法：}①改变金属组成（不锈钢）；②覆盖保护层（涂漆、电镀、搪瓷、钝化）；③电化学保护：\textbf{牺牲阳极法}（原电池原理，连接更活泼的 Zn、Mg 块作负极，被保护金属作正极；锌块会逐渐消耗，需定期更换）和\textbf{外加电流法}（电解原理，被保护金属接电源\textbf{负极}作阴极，辅助阳极接正极）。腐蚀快慢：电解池阳极$>$原电池负极$>$化学腐蚀$>$有防护措施的金属。

\begin{liti}[析氢腐蚀与吸氧腐蚀的实验判断]
将两枚相同的铁钉分别放入两支具支试管：甲中为食盐水浸过的铁钉，乙中为稀醋酸浸过的铁钉，支管均连接插入红墨水中的导管。一段时间后，甲导管中液面上升，乙导管中液面下降。解释原因，并写出甲中正极反应式。
\tcblower
\jie 甲中食盐水呈中性，铁发生吸氧腐蚀，消耗试管内的 \ce{O2}，气体压强减小，导管液面上升；正极：\ce{O2 + 4e- + 2H2O = 4OH-}。乙中稀醋酸呈酸性，发生析氢腐蚀，生成 \ce{H2} 使压强增大，导管液面下降。
\end{liti}

\begin{zhenti}{2026年广东卷第6题\quad 来源：仓库所附真题文档（网络收集版）}
做好防腐是发展海上风电的基础。下列措施\textbf{不能}达到防腐目的的是：A．钢制管桩连接铜块；B．钢制管桩外加富铝涂层；C．钢制管桩接外电源，另一端连接惰性电极；D．钢制管桩连接锌块。

\textbf{解析：}A 中 Cu 不如 Fe 活泼，Fe 作负极被加速腐蚀，不能防腐；B 富铝涂层既隔绝海水，铝又比铁活泼可作牺牲阳极；C 为外加电流法（管桩应接电源负极）；D 为牺牲阳极法。答案：A。
\end{zhenti}

% ---------------- 第5章 章末 ----------------
\zhangmo{本章小结}
\subsection*{一、知识总结}
\begin{center}
\small
\begin{tabularx}{\linewidth}{p{1.6cm}XX}
\toprule
& 原电池 & 电解池 \\
\midrule
能量转化 & 化学能$\to$电能（自发） & 电能$\to$化学能（非自发也能进行） \\
电极名称 & 负极（氧化）、正极（还原） & 阳极（氧化，接电源正极）、阴极（还原，接电源负极） \\
离子迁移 & 阳离子$\to$正极，阴离子$\to$负极 & 阳离子$\to$阴极，阴离子$\to$阳极 \\
典型应用 & 化学电源、加快反应、牺牲阳极法 & 氯碱、电镀、精炼、电冶金、外加电流法 \\
\bottomrule
\end{tabularx}
\end{center}

\subsection*{二、方法提升}
\begin{enumerate}
\item \textbf{判电极}：看化合价变化（升高为负极/阳极），或看外接电源，或看现象（溶解、析出、放气）。
\item \textbf{写反应}：先总后分、按介质补粒子；燃料电池“负极$=$总反应$-$正极”。
\item \textbf{判膜}：从电荷平衡和产品需求出发，确定哪种离子必须通过。
\item \textbf{算数量}：抓住电子守恒，建立 $e^-$ 与各物质的关系式。
\end{enumerate}

\subsection*{三、易错点归纳}
\begin{enumerate}
\item 认为活泼金属一定作负极（Mg–Al–NaOH、Fe–Cu–浓硝酸是反例）。
\item 碱性介质中电极反应出现 \ce{H+} 或 \ce{CO2}；酸性介质中出现 \ce{OH-}。
\item 认为电子可以通过电解质溶液；把离子迁移方向弄反。
\item 电解时忽略活性阳极自身溶解；电解后复原“加错物质”。
\item 外加电流法中把被保护金属接电源正极（应接负极）。
\end{enumerate}

\zhangmo{分层练习}
\zu{A组\quad 基础巩固}
\begin{xiti}
\item 关于 Zn–Cu 双液原电池（盐桥中为 KCl 琼胶），下列说法正确的是\xz
\xxa{电子由锌片经盐桥流向铜片}{盐桥中的 \ce{K+} 移向 \ce{ZnSO4} 溶液}{铜片上发生的反应为 \ce{Cu^2+ + 2e- = Cu}}{取出盐桥，电流计指针仍发生偏转}

\item 将镁片和铝片用导线相连插入 NaOH 溶液中，下列说法正确的是\xz
\xxa{镁比铝活泼，镁作负极}{负极反应为 \ce{Al - 3e- + 4OH- = [Al(OH)4]-}}{正极上产生 \ce{O2}}{电子由镁片经导线流向铝片}

\item 以 KOH 溶液为电解质的甲烷燃料电池中，负极反应式为\xz
\xxa{\ce{CH4 - 8e- + 2H2O = CO2 + 8H+}}{\ce{CH4 - 8e- + 10OH- = CO3^2- + 7H2O}}{\ce{O2 + 4e- + 2H2O = 4OH-}}{\ce{CH4 + 8e- + 10OH- = CO3^2- + 7H2O}}

\item 关于铅酸蓄电池，下列说法正确的是\xz
\xxa{放电时，负极质量减小}{放电时，电解质溶液的密度增大}{充电时，铅酸蓄电池的负极应与外接电源的正极相连}{充电时，阳极反应为 \ce{PbSO4 - 2e- + 2H2O = PbO2 + 4H+ + SO4^2-}}

\item 用惰性电极电解足量 \ce{CuSO4} 溶液一段时间后，要使溶液恢复到原来的浓度和体积，应加入适量的\xz
\xx{Cu}{CuO}{\ce{Cu(OH)2}}{\ce{CuSO4}}

\item \laiyuan{2024年广东卷第5题，来源：仓库所附真题文档}\ 我国自主设计建造的浮式生产储卸油装置“海葵一号”将在珠江口盆地海域使用，其钢铁外壳镶嵌了锌块，以利用电化学原理延缓外壳的腐蚀。下列有关说法正确的是\xz
\xxb{钢铁外壳为负极}{镶嵌的锌块可永久使用}{该法为外加电流法}{锌发生反应：\ce{Zn - 2e- = Zn^2+}}

\item 用铜作电极电解稀硫酸，下列说法正确的是\xz
\xxa{阳极上放出 \ce{O2}}{阴极上析出 Cu}{溶液逐渐变蓝，阴极上放出 \ce{H2}}{电解一段时间后溶液 pH 减小}
\end{xiti}

\zu{B组\quad 高考典型}
\begin{xiti}
\item \laiyuan{2025年广东卷第16题，来源：仓库所附真题文档}\ 某理论研究认为：氢氧燃料电池（图b，两电极间为质子交换膜）的电极Ⅰ和Ⅱ上所发生反应的催化机理分别如图a和图c。图a：\ce{O2} 在电极Ⅰ表面的催化剂上依次经过 i$\to$ii$\to$iii$\to$iv 四步，每步结合 1 个 \ce{H+} 和 1 个 $e^-$，最终生成 \ce{H2O}；图c：\ce{H2} 在电极Ⅱ表面经两步各失去 1 个电子生成 \ce{H+}。其中 \ce{O2} 获得第一个电子的过程（i$\to$ii）最慢。由此可知，理论上\xz
\xxb{负极反应的催化剂是 i}{图a中，i 到 ii 过程的活化能一定最低}{电池工作过程中，负极室的溶液质量保持不变}{相同时间内，电极Ⅰ和电极Ⅱ上的催化循环完成次数相同}

\item \laiyuan{原创}\ 磷酸铁锂电池的总反应为 \ce{Li_xC6 + Li_{1-$x$}FePO4 <=>[\text{放电}][\text{充电}] 6C + LiFePO4}，电解质为能传导 \ce{Li+} 的有机溶液。(1) 写出放电时正极的电极反应式。(2) 充电时，\ce{Li+} 向哪一极迁移？阳极反应式是什么？(3) 放电时每转移 1 mol 电子，负极材料质量变化多少？(4) 为什么该电池不能用水溶液作电解质？

\item 下列关于电镀和电解精炼的说法正确的是\xz
\xxa{在铁件上镀锌，铁件作阳极，锌片作阴极}{电镀过程中，电解质溶液中镀层金属离子的浓度明显减小}{电解精炼铜时，粗铜作阳极，阳极减少的质量等于阴极增加的质量}{电解精炼铜时，Ag、Au 等不如铜活泼的金属沉积在阳极底部形成阳极泥}

\item \laiyuan{原创}\ 用如下三室装置电解 \ce{Na2SO4} 溶液制备 NaOH 和 \ce{H2SO4}：左侧为阳极室（稀硫酸），中间室为 \ce{Na2SO4} 溶液，右侧为阴极室（稀 NaOH 溶液），两张膜分别为 X 和 Y（X 在阳极室与中间室之间）。(1) X、Y 分别是阴离子交换膜还是阳离子交换膜？(2) 写出阳极、阴极的电极反应式。(3) 以 2.0 A 电流电解 965 s，理论上阴极室生成 NaOH 的物质的量为多少？中间室 \ce{Na2SO4} 减少多少克？

\item \laiyuan{原创}\ 锌—空气电池以 KOH 溶液为电解质，总反应为 \ce{2Zn + O2 + 4OH- + 2H2O = 2[Zn(OH)4]^2-}。(1) 写出正、负极的电极反应式。(2) 放电时，\ce{OH-} 向哪一极移动？电解质溶液中 $c(\ce{OH-})$ 如何变化？(3) 每消耗 6.5 g Zn，理论上消耗标准状况下空气多少升（空气中 \ce{O2} 体积分数按 20\% 计）？

\item \laiyuan{原创}\ 将两个电解池串联：甲池为惰性电极电解足量 \ce{AgNO3} 溶液，乙池为电解精炼铜（粗铜含少量 Zn、Ag）。通电一段时间后甲池阴极增重 10.8 g。(1) 乙池阴极增重多少克？(2) 乙池阳极减少的质量是否等于阴极增加的质量？为什么？(3) 甲池阳极产生气体在标准状况下的体积为多少？
\end{xiti}

\zu{C组\quad 冲刺高分}
\begin{xiti}
\item \laiyuan{2026年广东卷第16题，来源：仓库所附真题文档（网络收集版）}\ 一种利用电解还原 \ce{CO2} 同时实现海水脱盐的装置：从左到右依次为电极a、Ⅰ室、膜①、Ⅱ室（海水）、膜②、Ⅲ室（海水）、膜③、Ⅳ室、电极b；Ⅰ室中通入 \ce{CO2} 并放出 CO，Ⅳ室放出 \ce{O2}；电解液在Ⅰ和Ⅳ室间循环，初始电解液为 \ce{NaHCO3} 溶液；膜①～③为阳离子或阴离子交换膜。下列说法错误的是\xz
\xxb{①为阳离子交换膜，②为阴离子交换膜}{阳极反应为 \ce{4OH- - 4e- = 2H2O + O2 ^}}{电解一段时间后，Ⅲ室海水中盐的浓度降低}{电解过程中，Ⅰ室电解液的 pH 在一段时间内可保持稳定}

\item \laiyuan{原创}\ 以质子交换膜甲醇燃料电池为电源，用惰性电极电解 200 mL 0.50 mol·L$^{-1}$ \ce{CuSO4} 溶液。(1) 写出燃料电池负极、正极的电极反应式。(2) 当燃料电池消耗 0.32 g 甲醇时，电解池阴极析出铜多少克？阳极产生气体在标准状况下的体积为多少？(3) 此时电解液中 $c(\ce{H+})$ 为多少（忽略体积变化）？要使电解液复原，应加入什么物质、多少克？(4) 若继续电解至阴极开始产生气体，此后电解的实质是什么？

\item \laiyuan{原创}\ 用电解法将 \ce{Na2CrO4} 转化为 \ce{Na2Cr2O7}：电解槽用阳离子交换膜隔开，阳极室加入 \ce{Na2CrO4} 溶液，阴极室加入稀 NaOH 溶液，两极均为惰性电极。已知 \ce{2CrO4^2- + 2H+ <=> Cr2O7^2- + H2O}。(1) 写出阳极、阴极的电极反应式。(2) 解释阳极室中 \ce{Na2CrO4} 转化为 \ce{Na2Cr2O7} 的原理。(3) 通过阳离子交换膜迁移的离子是什么？向哪个方向迁移？(4) 若阳极室溶液中 Na 与 Cr 的物质的量之比由开始的 2:1 变为 $d$:1，则此时 \ce{Na2CrO4} 的转化率为\tk（用 $d$ 表示）。
\end{xiti}

\zhangmo{练习解析}
\begin{jiexilist}
\item \ans{C}电子只经导线流动，A 错；盐桥中阳离子 \ce{K+} 移向正极区（\ce{CuSO4} 溶液），B 错；铜为正极，\ce{Cu^2+} 得电子，C 对；取出盐桥，回路断开，无电流，D 错。
\item \ans{B}Mg 不与 NaOH 反应，Al 与 NaOH 反应，Al 作负极，A 错、B 对；正极上 \ce{H2O} 得电子生成 \ce{H2}，C 错；电子由负极（铝）流向正极（镁），D 错。
\item \ans{B}碱性介质中 C 元素以 \ce{CO3^2-} 形式存在，不能写 \ce{H+} 和 \ce{CO2}；C 是正极反应；D 中负极应失电子。
\item \ans{D}放电时负极 Pb 生成 \ce{PbSO4}，质量增大，A 错；消耗 \ce{H2SO4}、生成水，溶液密度减小，B 错；充电时“负接负”，C 错；充电阳极是放电正极反应的逆反应，D 对。
\item \ans{B}电解时阴极析出 Cu，阳极放出 \ce{O2}，“出什么加什么”，相当于减少了 CuO，应加 CuO；加 \ce{Cu(OH)2} 会多加水。
\item \ans{D}Zn 比 Fe 活泼，锌块为负极，钢铁外壳为正极被保护，A 错；锌不断溶解，需定期更换，B 错；该法为牺牲阳极法，C 错；D 对。
\item \ans{C}阳极 Cu 为活性电极，自身溶解：\ce{Cu - 2e- = Cu^2+}，溶液变蓝；阴极开始时 \ce{H+} 放电产生 \ce{H2}（\ce{H+} 浓度大，生成的 \ce{Cu^2+} 较少），总反应 \ce{Cu + H2SO4 ->[\text{电解}] CuSO4 + H2 ^}，消耗 \ce{H+}，pH 增大。
\item \ans{C}电极Ⅰ上 \ce{O2} 被还原，为正极，i 是正极反应的催化剂，A 错；i$\to$ii 最慢，活化能最高，B 错；负极 \ce{H2 - 2e- = 2H+}，每失去 1 个电子就有 1 个 \ce{H+} 通过质子交换膜进入正极室，负极室溶液质量不变，C 对；\ce{O2} 完成一次循环需 4 个电子，\ce{H2} 完成一次循环需 2 个电子，同一时间转移电子数相同，两者循环次数之比为 1:2，D 错。
\item (1) \ce{Li_{1-$x$}FePO4 + $x$Li+ + $x$e- = LiFePO4}。(2) 充电时阳离子移向阴极（即放电时的负极，石墨一侧）；阳极：\ce{LiFePO4 - $x$e- = Li_{1-$x$}FePO4 + $x$Li+}。(3) 负极 \ce{Li_xC6 - $x$e- = $x$Li+ + 6C}，每转移 1 mol 电子，负极减少 1 mol Li，即质量减少 7 g。(4) Li 及嵌锂石墨很活泼，会与水剧烈反应。
\item \ans{D}镀锌时铁件作阴极、锌片作阳极，A 错；电镀时阳极溶解的金属离子与阴极析出的相等，浓度基本不变，B 错；精炼时阳极上 Zn、Fe、Ni 等也溶解且有阳极泥，两极质量变化不相等，C 错；D 对。
\item (1) 阳极室需要 \ce{SO4^2-} 进入（与生成的 \ce{H+} 组成 \ce{H2SO4}），X 为阴离子交换膜；阴极室需要 \ce{Na+} 进入（与生成的 \ce{OH-} 组成 NaOH），Y 为阳离子交换膜。(2) 阳极：\ce{2H2O - 4e- = O2 ^ + 4H+}；阴极：\ce{2H2O + 2e- = H2 ^ + 2OH-}。(3) $n(e^-)=\dfrac{2.0\times965}{96500}=0.020$ mol；生成 NaOH 0.020 mol；中间室迁出 0.020 mol \ce{Na+} 和 0.010 mol \ce{SO4^2-}，即减少 0.010 mol \ce{Na2SO4}，1.42 g。
\item (1) 负极：\ce{Zn - 2e- + 4OH- = [Zn(OH)4]^2-}；正极：\ce{O2 + 4e- + 2H2O = 4OH-}。(2) 移向负极；总反应消耗 \ce{OH-}，$c(\ce{OH-})$ 减小。(3) 0.1 mol Zn 失 0.2 mol 电子，需 0.05 mol \ce{O2}，即 1.12 L，空气 $1.12\div20\%=5.6$ L。
\item (1) 甲池析出 Ag 0.1 mol，转移 0.1 mol 电子；乙池阴极析出 Cu 0.05 mol，增重 3.2 g。(2) 不相等。粗铜中的 Zn 也失电子溶解（摩尔质量与 Cu 不同），Ag 不溶解而落入阳极泥，所以阳极减少的质量与阴极增加的质量不相等，但两极转移的电子数相等。(3) 甲池阳极：\ce{2H2O - 4e- = O2 ^ + 4H+}，0.1 mol 电子生成 0.025 mol \ce{O2}，即 0.56 L。
\item \ans{C}a 极 \ce{CO2} 得电子生成 CO，为阴极：\ce{CO2 + H2O + 2e- = CO + 2OH-}，生成的 \ce{OH-} 与 \ce{HCO3-} 反应生成 \ce{CO3^2-}；b 极为阳极，水电离的 \ce{OH-} 放电生成 \ce{O2}，B 对。若膜③为阴离子交换膜，Ⅲ室 \ce{Cl-} 会进入Ⅳ室在阳极放电生成 \ce{Cl2}，故③应为阳离子交换膜（Ⅳ室 \ce{Na+} 进入Ⅲ室）；要使某室脱盐，该室的 \ce{Na+}、\ce{Cl-} 必须都能移出：Ⅱ室 \ce{Na+} 经膜①（阳膜）移入Ⅰ室，\ce{Cl-} 经膜②（阴膜）移入Ⅲ室，A 对。于是Ⅱ室盐浓度降低（脱盐），Ⅲ室得到 \ce{Cl-}（来自Ⅱ）和 \ce{Na+}（来自Ⅳ），盐浓度升高，C 错。Ⅰ室中 \ce{HCO3-}/\ce{CO3^2-} 构成缓冲体系，且与Ⅳ室电解液循环，pH 可在一段时间内保持稳定，D 对。
\item (1) 负极：\ce{CH3OH - 6e- + H2O = CO2 ^ + 6H+}；正极：\ce{O2 + 4e- + 4H+ = 2H2O}。(2) $n(\ce{CH3OH})=0.01$ mol，转移 0.06 mol 电子。$n(\ce{Cu^2+})=0.10$ mol，足量，析出 Cu 0.03 mol，即 1.92 g；阳极 \ce{O2} 0.015 mol，0.336 L。(3) 生成 \ce{H+} 0.06 mol，$c(\ce{H+})=0.30$ mol·L$^{-1}$；溶液中减少的是 0.03 mol Cu 和 0.015 mol \ce{O2}，相当于 0.03 mol CuO，加入 CuO 2.4 g（或 \ce{CuCO3} 3.72 g）。(4) \ce{Cu^2+} 完全放电后，阴极上 \ce{H+}（来自水和生成的 \ce{H2SO4}）放电，阳极仍放出 \ce{O2}，实质是电解水。
\item (1) 阳极：\ce{2H2O - 4e- = O2 ^ + 4H+}；阴极：\ce{2H2O + 2e- = H2 ^ + 2OH-}。(2) 阳极生成 \ce{H+}，使 \ce{2CrO4^2- + 2H+ <=> Cr2O7^2- + H2O} 正向移动。(3) \ce{Na+}，由阳极室迁移到阴极室（阴极室得到 NaOH）。(4) $2-d$。设起始 \ce{Na2CrO4} 1 mol（Na 2 mol，Cr 1 mol），转化率为 $\alpha$：生成 $\alpha$ mol \ce{H+} 需转移 $\alpha$ mol 电子，同时有 $\alpha$ mol \ce{Na+} 迁出阳极室，剩余 Na $(2-\alpha)$ mol，Cr 仍为 1 mol。$\dfrac{2-\alpha}{1}=d$，$\alpha=2-d$。
\end{jiexilist}

